% !PRINTED 52
\weberpage{52}

%<<CONT
or else, if one transforms $u$ into $u+2iK'$ and again drops the notation $\varkappa$:
\weq{13}{\mathrm{sn}\,(u+2iK') = \mathrm{sn}\,u.}
The function sn $u$ therefore has a \gesperrt{double periodicity}. One of these periods, $4K$, is real; the other, $2iK$, purely imaginary % [sic] German source prints K without the mathematically expected prime.
 (at least when the modulus $\varkappa$ is a positive proper fraction). This double periodicity is a property which occurs nowhere in the region of the elementary functions, and which thus showed most clearly that one has here to do with a new kind of function. At the same time it is the consequences derived from the double periodicity which, as experience has shown, lead by far the most quickly and with satisfying rigour into the interior of the theory, so that the most suitable starting point for a methodical development is to be sought here. Our next considerations will therefore be devoted to the doubly periodic functions in general.
\begin{center}
\rule{3cm}{0.4pt}
\end{center}
